\documentclass{aevidence}

\journalvolume{14}
\journalissue{2}
\journalyear{2026}
\articledate{27.11.2026}
\articlepages{35--51}

\title{Empirische Plausibilität bei unvollständiger Datenlage}

\subtitle{Eine statistische Untersuchung zur Frage, wann ein Ergebnis überzeugend genug klingt}

\author{
  Dr. Konrad P. Huber
  \and
  Dr. Maximilian Methodius
}

\shorttitlecommand{Empirische Plausibilität}

\articlekeywords{
  Plausibilität; Datenlage; empirische Evidenz; Wahrnehmung
}

\abstracttext{
Diese Untersuchung befasst sich mit der Frage, unter welchen Bedingungen
empirische Ergebnisse auch bei unvollständiger Datenlage als hinreichend
plausibel wahrgenommen werden. Dabei wird insbesondere untersucht, welchen
Einfluss die sprachliche Darstellung eines Ergebnisses auf dessen
wahrgenommene Überzeugungskraft besitzt.
}

\begin{document}

\section{Einleitung}

Empirische Untersuchungen setzen üblicherweise eine hinreichende Datenbasis
voraus. In der wissenschaftlichen Praxis zeigt sich jedoch, dass Ergebnisse
auch dann als überzeugend wahrgenommen werden können, wenn die zugrunde
liegende Datenlage nur teilweise vorliegt.

Die vorliegende Untersuchung widmet sich daher der Frage, ob die
wahrgenommene Plausibilität eines Ergebnisses in einem systematischen
Zusammenhang mit dem Umfang der tatsächlich verfügbaren Daten steht.

\section{Methodik}

Für die Untersuchung wurden verschiedene hypothetische Datensätze
konstruiert. Anschließend wurden daraus jeweils kurze wissenschaftliche
Aussagen formuliert.

Besonderes Augenmerk wurde auf die sprachliche Sicherheit der Aussagen
gelegt. Formulierungen wie „die Ergebnisse zeigen“ wurden dabei solchen
gegenübergestellt, die auf eine geringere Sicherheit hinweisen.

\section{Ergebnisse}

Die Ergebnisse zeigen einen deutlichen Zusammenhang zwischen der
sprachlichen Sicherheit einer Aussage und ihrer wahrgenommenen empirischen
Belastbarkeit.

Bemerkenswert war insbesondere, dass die Beurteilung der Plausibilität nur
in begrenztem Umfang von der tatsächlich verfügbaren Datenmenge beeinflusst
wurde.

\section{Diskussion}

Die Befunde legen nahe, dass die Überzeugungskraft eines Ergebnisses nicht
ausschließlich von der Qualität der zugrunde liegenden Daten abhängt.

Eine abschließende Bewertung dieses Zusammenhangs setzt allerdings weitere
Untersuchungen voraus, insbesondere solche, bei denen die Datenlage
vollständig bekannt ist.

\section{Fazit}

Ein Ergebnis kann auch dann plausibel erscheinen, wenn die empirische
Grundlage nicht vollständig vorliegt.

Ob dies als methodisches Problem oder als effiziente Form der
wissenschaftlichen Kommunikation zu bewerten ist, bleibt Gegenstand
weiterer Untersuchungen.

\end{document}
